\begin{apartats}

\apartat[1,25]{0,75}
Estudia la posició relativa de la recta
$r\colon\left\{\begin{aligned}x&=1+2t\\y&=-1+t\\z&=3-t\end{aligned}\right.$ i el pla
$\pi\colon x-y+z-1=0$.

\begin{solucio}
Substituint la recta al pla: $(1+2t)-(-1+t)+(3-t)-1=4$, i l'equació $4=0$ no té solució. La recta i el
pla no tenen cap punt en comú: són paral·lels. (Ho confirma que $\vec u=(2,1,-1)$ i $\vec n=(1,-1,1)$
compleixen $\vec u\cdot\vec n=2-1-1=0$.)
\end{solucio}

\begin{tria}{recta-pla}
\itemtria{parametres}{1}{1,25}
Estudia, segons els valors de $a$ i $b$, la posició relativa del pla $ax+2y-z+b=0$ i la recta
\[
\frac{x-1}{1}=\frac{y+1}{2}=\frac{z-3}{-1}.
\]

\begin{solucio}
En paramètriques, la recta és $(1+t,\ -1+2t,\ 3-t)$. Substituint:
$a(1+t)+2(-1+2t)-(3-t)+b=(a+5)t+(a+b-5)=0$.
\begin{itemize}
\item Si $a\ne-5$, hi ha un sol valor de $t$: es tallen en un punt.
\item Si $a=-5$, queda $b-10=0$. Si $b=10$, qualsevol $t$ la compleix: la recta és continguda al pla.
Si $b\ne10$, cap: són paral·lels.
\end{itemize}
\end{solucio}

\itemtria{punt-de-tall}{1}{1,25}
Comprova que la recta
$s\colon\left\{\begin{aligned}x&=2+t\\y&=-t\\z&=1+2t\end{aligned}\right.$ talla el pla
$x+y+z-7=0$, i troba el punt de tall.

\begin{solucio}
Substituint: $(2+t)+(-t)+(1+2t)-7=2t-4=0$, $t=2$. Hi ha una sola solució: es tallen al punt
$(4,-2,5)$.
\end{solucio}
\end{tria}

\begin{nomesllarg}
\apartat{0,75}
Considera els plans $\pi_1\colon ax+2y-4z=1$ i $\pi_2\colon x+y-2z=3$. Troba el valor de $a$ perquè
siguin paral·lels, i el valor de $a$ perquè siguin perpendiculars.

\begin{solucio}
Paral·lels: els vectors normals $(a,2,-4)$ i $(1,1,-2)$ han de ser proporcionals, $\frac a1=\frac21=\frac{-4}{-2}$,
és a dir, $a=2$. Com que $\frac13\ne2$, no són coincidents.\\
Perpendiculars: els vectors normals han de ser perpendiculars, $a+2+8=0$, és a dir, $a=-10$.
\end{solucio}
\end{nomesllarg}

\end{apartats}
